\subsection{The groups $C_2\times C_{36}$ and $C_{72}$}
\label{subsec:decic236}

Put $X=X_1(2,36)$. This curve has genus $41$ and good reduction
at $5$. Its Jacobian has rank zero over $\mathbb Q$ by
\cite[Theorem~3.1(3)]{DEHMZ}.
We first exclude degree-$10$ divisors whose reduction has a
noncuspidal component. We then exclude degree-$10$ pencils and
use torsion reduction to deal with entirely cuspidal reduction.
All effective divisors considered below may be reducible.

\begin{lemma}\label{lem:decic236-quotient-image}
Let $X=X_1(2,36)$ and $h(P,Q)=(P,13Q)$, and put
$Z=X/\langle h\rangle$. The curve $Z$ has genus $9$.
The image of $J_Z(\mathbb Q)$ in $J_Z(\mathbb F_5)$ is contained
in the subgroup generated by the finite cuspidal classes. This
subgroup has invariants $[3,3,3,24,24,24]$ and index $8$.
\end{lemma}
\begin{proof}
Here is an explicit model of the quotient. Over $B=X_1(2,12)$
write $(E,P,Q_{12})$ for the universal level structure, put
$T=4Q_{12}$, and let $E'=E/\langle T\rangle$.
The $3$-division polynomial of $E'$ has a linear factor for the
dual kernel. Adjoin a root of its remaining cubic factor.
This gives $Z$.
Indeed, the nine lifts $Q_{36}$ with $3Q_{36}=Q_{12}$ split into
three orbits under $h$. Their images $\phi(4Q_{36})$ give the three cyclic subgroups of
order $3$ distinct from the dual kernel.
The choice of their $x$-coordinate gives the stated cubic factor.
The model has genus $9$ and good reduction at $5$ and $7$.
Its Jacobian has rank zero, since it is a quotient of $J_X$.

Put $A_7=T_7-1-7\langle7\rangle$.
The correspondence descends from $X$, and Eichler--Shimura and
odd unramified torsion injectivity at $7$ give
$A_7J_Z(\mathbb Q)=0$.
At $5$ the proven finite Jacobian is
\[
 G=[3,3,3,24,24,192].
\]
Let $H$ be its subgroup generated by all finite cuspidal classes.
The exact computation gives $H=[3,3,3,24,24,24]$ and $G/H=C_8$.

We determine the action on these cusps by Tate labels.
For $Q_{36}=\zeta_{36}^a q^{b/36}$ and
$P=(-1)^c q^{d/2}$, impose $cb-da\equiv1\pmod2$ and quotient
by the unipotent cusp action and simultaneous negation.
The $64$ source cusps give $32$ cusps on $Z$.
Their global closed degrees are $1^4,2^{10},4^2$; at $5$ they
split as $1^8,2^{12}$.
Thirteen rational modular-unit valuation signatures distinguish
all sixteen global orbits. They identify the finite Galois
permutation $\sigma_7$ and diamond permutation $\delta_7$
unambiguously, including the four split pairs.
The boundary formula
\[
 T_7c=\sigma_7c+7\delta_7\sigma_7^{-1}c
\]
then gives $A_7H=0$ as classes.

A noncuspidal cubic place $z$ gives a class
$a=[z-3c_0]$ generating $G/H$, where $c_0$ is a global rational
cusp. Exact Riemann--Roch calculations show that $A_7a$ has order
$8$: its multiples $1,2,4$ are nonprincipal and its multiple $8$
is principal with an exported defining function.
The projected Hecke divisor has degree $24$; the kernel-orbit
algorithm and the untouched full-division-field algorithm give
identical divisors. The assignments on the twenty cusp generators
and $a$ respect every relation of $G$, independently checked by
integer lattice arithmetic. Consequently
$\ker(A_7:G\to G)=H$.
Reduction at $5$ proves the assertion.
\end{proof}

Let $D$ be an effective divisor on $X$, defined over $\mathbb Q$,
of arbitrary degree. Write its reduction at $5$ as $N+C$, with
$N$ noncuspidal and $C$ cuspidal, and let $\pi:X\to Z$ be the
quotient map. The normalized class
$[\pi_*D-(\deg D)c_0]$ lies in $J_Z(\mathbb Q)$, and hence reduces
to $H$. The normalized cuspidal class
$[\pi_*C-(\deg C)c_0]$ also lies in $H$.
It follows that
\[
 [\pi_*N-(\deg N)c_0]\in H.
\]
This necessary condition applies, in particular, when $\deg D=10$.
For a closed point of degree $d$ whose image has degree $e$, its
pushforward is $d/e$ times that image point; we include these
multiplicities in every class calculation.


\begin{lemma}\label{lem:decic236-finite-sieve}
Let $D$ be an effective divisor of degree $10$ on $X$, defined
over $\mathbb Q$, whose support is entirely noncuspidal.
Suppose that its reduction at $5$ has a noncuspidal component.
After applying a rational level automorphism, its reduction is
$\overline N+\overline C$, where $\overline N$ is the degree-$8$
point in Lemma~\ref{lem:236-degree10-oldN8} and $\overline C$
is a cuspidal effective divisor of degree $2$.
\end{lemma}
\begin{proof}
The exact noncuspidal closed-point counts over $\mathbb F_5$ in
degrees one through ten are
\[
 (0,0,48,120,720,2904,10416,50334,218080,979512).
\]
The cusp signature is $1^{24},2^8,6^4$.
The point counts agree with Newton identities for the independently
computed Frobenius polynomial.
The $24$ rational level automorphisms are
\[
 (P,Q)\longmapsto(P+18bQ,cP+aQ),\qquad
 b,c\in\{0,1\},\quad a\in\{1,5,7,11,13,17\}.
\]
Their noncuspidal orbit counts in degrees three through ten are
\[
 (2,5,30,126,434,2115,9088,40874).
\]

Write the reduction as $N+C$, where $N$ is noncuspidal and
$C$ is cuspidal. Since the smallest noncuspidal closed degree is
$3$, the divisor $N$ has one, two, or three summands, including
repetitions. We normalize one summand by the level group and
retain every actual closed point for the remaining summands.
This gives $135282$ support patterns and includes every possible
$N$ with $3\leq\deg N\leq10$.

First, we use two maps to $Y=X_1(2,18)$:
\[
 (E,P,Q)\longmapsto(E,P,2Q),\qquad
 (E,P,Q)\longmapsto
 (E/\langle18Q\rangle,\phi(P),\phi(Q)).
\]
The rational Jacobian of $Y$ is the cuspidal group
$[2,42,126]$ by \cite[Proposition~4.16 and Table~3]{DEHMZ}.
Its reduced image has index $12$ in
$J_Y(\mathbb F_5)=[6,84,252]$.
Every normalized finite cusp class lies in this image.
Thus the two normalized pushforward classes of $N$ must lie in
this subgroup. These conditions leave $1251$ patterns.

Next, we use the genus-$9$ quotient and
Lemma~\ref{lem:decic236-quotient-image}.
Its image condition leaves $208$ patterns.
Finally, we use four maps to $V=X_1(36)$.
The first two retain $Q$ or $P+Q$.
The other two retain the image of $Q$ after quotienting by
$\langle P\rangle$ or $\langle P+18Q\rangle$.
The quotient kernels are disjoint from $\langle Q\rangle$, so
the resulting points have order $36$.

By \cite[Theorem~4.13 and Table~2]{DEHMZ}, the rational
Jacobian of $V$ is the cuspidal group $[12,252,7812]$.
The proven finite Jacobian and reduced rational image are
\[
 G=[4,4,12,12,36,7812,7812],\qquad H=[12,252,7812],
\]
with $[G:H]=214272$.
Let $q:G\to G/H$ be the quotient map, and fix a rational cusp
$c_0$ on $V$. For $0\leq e\leq7$, let $\mathcal C_e$ be the
set of classes $q([A-ec_0])$, where $A$ ranges over finite
cuspidal effective divisors of degree $e$.
Their exact cardinalities are
\[
 (\#\mathcal C_e)_{e=0}^7=(1,7,27,77,181,273,279,279).
\]
For each of the four maps $f$, the necessary condition is
\[
 -q\bigl([f_*N-(\deg N)c_0]\bigr)\in\mathcal C_{10-\deg N}.
\]
We include the norm multiplicity $d/e$ whenever a closed point of
degree $d$ maps to one of degree $e$.

These four tests leave precisely one pattern: the point
$\overline N$ of degree $8$ in
Lemma~\ref{lem:236-degree10-oldN8}, with a cuspidal complement
of degree $2$. The field modulus and both coordinate vectors
agree exactly with those in that lemma.

Every class group is computed with proof.
Independent finite-field arithmetic checks all projection values,
residue degrees, and class sums.
A second native calculation checks the finite and infinite ideal
identities for all $57384$ genus-$7$ target images,
all $1104$ genus-$9$ target images, and all $240$ final target
images. Of these, respectively $16$, $82$, and $129$ lie directly
in the proven factorbase; all others have exact exported
Riemann--Roch identity witnesses.
This proves the assertion.
\end{proof}

% Fragment for the degree-ten C_2 x C_36 exclusion. No irreducibility is assumed.
\begin{lemma}\label{lem:236-degree10-oldN8}
Let $X=X_1(2,36)$, and let $\overline N$ be the noncuspidal
closed point of degree $8$ over $\mathbb F_5$ specified below.
There is no effective divisor $D$ on $X$, defined over $\mathbb Q$,
of degree $10$, whose support is entirely noncuspidal and whose
reduction at $5$ is $\overline N+\overline C$, where
$\overline C$ is a cuspidal effective divisor of degree $2$.
\end{lemma}
\begin{proof}
Put $g(P,Q)=(P+18Q,P+Q)$.
This automorphism has order $4$ and acts freely on the geometric
cusps. The point $\overline N$ is invariant under $g$.
Write $\pi:X\longrightarrow W=X/\langle g\rangle$.
The curve $W$ has genus $11$, the map $\pi$ is unramified, and
$\overline N=\pi^*q$, where $q$ has degree $2$ over $\mathbb F_5$.

We first compute
\begin{equation}\label{eq:236-degree10-cusp-section}
 h^0\bigl(\overline N+\pi^*\pi_*\overline C\bigr)=1
\end{equation}
for every cuspidal effective divisor $\overline C$ of degree $2$.
Here is a model of $W$ and the required calculation.
Put $r=t^2-t+1$, $c=t(r-1)$, $b=rc$, and
\[
 E_9:\quad y^2+(1-c)xy-by=x^3-bx^2.
\]
Then $W$ is given by
\[
 16(w^2+16)^3=j(E_9)w^4.
\]
After completing the square in $E_9$, write its equation as
$y^2=x^3+a_2x^2+a_4x+a_6$.
The cubic has a unique common root $\rho$ with
\[
 w^2(a_2+3x)^2-4(w^2+4)(a_4+2a_2x+3x^2).
\]
Put $\beta=a_4+2a_2\rho+3\rho^2$ and
$f=16\beta^3(1+4/w)^2$.
Over $\mathbb F_5$ the cover is $z^4=f$.
Its divisor is $\operatorname{div}(f)=4D_0$, where $[D_0]$ has
order $4$. Hence, for every divisor $A$ on $W$,
\[
 h^0_X(\pi^*A)=\sum_{j=0}^3h^0_W(A+jD_0).
\]

There are six rational cusps $c_1,\ldots,c_6$ of $W$ whose
fibres split into four rational cusps of $X$.
There are two quadratic cusps $u_1,u_2$ whose fibres split into
four quadratic cusps. These account for all source cusps of
degree at most $2$.
The exact Riemann--Roch calculations give
\[
 \bigl(h^0(A+jD_0)\bigr)_{j=0}^3=(1,0,0,0)
\]
for the two divisors $A=q+u_i$ and the $21$ divisors
$A=q+c_i+c_j$, where $1\leq i\leq j\leq6$.
This proves \eqref{eq:236-degree10-cusp-section}.

Suppose that $D$ exists. Work over a finite unramified extension
of $\mathbb Q_5$ containing the cusps, and let $C$ be the cuspidal
lift of $\overline C$. Put $L=[D-C]\in\operatorname{Pic}^8(X)$.
The Jacobian $J_X(\mathbb Q)$ has rank zero. Together with
Manin--Drinfeld, this implies that $gL-L$ is torsion.
Its reduction is zero, since $\overline N$ is $g$-invariant.
Torsion reduction is injective over an unramified extension of
$\mathbb Q_5$, and hence $gL=L$.

Put $\Omega=\pi^*\pi_*C=\sum_{j=0}^3g^jC$.
The divisors
\[
 D+\Omega-C\quad\hbox{and}\quad gD+\Omega-gC
\]
are effective and linearly equivalent.
Their common line bundle reduces to
$\mathcal O(\overline N+\pi^*\pi_*\overline C)$.
By \eqref{eq:236-degree10-cusp-section} and upper semicontinuity,
its space of sections has dimension at most $1$.
Thus the two divisors are equal.
Since $D$ and $gD$ have no cuspidal components, comparison of
cuspidal parts gives $C=gC$.
Every geometric cusp orbit has length $4$, whereas $\deg C=2$.
This is a contradiction.
\end{proof}

\noindent\emph{Identification and computational evidence.}
For an exact specification of $\overline N$, let
\[
 k=\mathbb F_5[\alpha]/
 (\alpha^8+\alpha^4+3\alpha^2+4\alpha+2).
\]
Write $B=\sum_{i=0}^7B_i\alpha^i$ and
$C=\sum_{i=0}^7C_i\alpha^i$, with coefficient vectors
\[
 (B_i)=(0,0,0,0,2,3,0,4),\qquad
 (C_i)=(2,1,0,4,0,2,1,2).
\]
The level structure
$E:y^2=x^3+Cx^2+B(1-B-C)x$, $P=(0,0)$, $Q=(B,B)$
represents $\overline N$. On it, $g$ agrees with the sixth
power of the $5$-power Frobenius.
The function-field calculation and all $23$ Riemann--Roch tests
are in \texttt{new236\_geom\_cusp2\_numeric.m}; the exported data
are in \texttt{geometry/cusp2\_numeric\_certificate.m}.
An independent exact replay in Sage verifies the level structure,
the Frobenius relation, the source-to-quotient invariants, the
local Kummer splitting, and the free action on all $64$ geometric
cusps; see \texttt{geometry/replay\_cusp2\_v2.py} and
\texttt{geometry/cusp2\_independent\_replay.json}.


\begin{lemma}\label{lem:decic236-cuspidal-pencil}
Every degree-$10$ morphism
$X_1(2,36)_{\mathbb F_5}\to\mathbb P^1$
has a rational fibre whose support is not entirely cuspidal.
\end{lemma}
\begin{proof}
The source has $24$ rational points, all cuspidal, and its cusp
signature is $1^{24},2^8,6^4$. Suppose that all six rational
fibres are cuspidal. Let $n_1\geq\cdots\geq n_6$ count the
distinct rational cusps in these fibres. Then
\[
 \sum_i n_i=24,\qquad n_1\geq4,\qquad
 n_2\geq\left\lceil\frac{24-n_1}{5}\right\rceil.
\]
Let $A_1,A_2$ be the two fibres with largest rational support.
For a cuspidal divisor $A=\sum_c m_c c$, use the geometric
coefficient norm
\[
 \|A\|^2=\sum_c \deg(c)m_c^2.
\]
If $\deg A=10$ and it contains $n$ distinct rational cusps,
then $\|A\|^2\leq(11-n)^2+n-1$.
The fibres have disjoint support. Checking $4\leq n_1\leq10$
and the displayed lower bound on $n_2$ gives
\[
 \|A_1-A_2\|^2\leq104.
\]
The positive and negative parts both have degree $10$.

Lift the cusps over a cyclotomic extension unramified at $5$.
Their difference is torsion by Manin--Drinfeld and reduces to a
principal class. Torsion injectivity over this unramified odd local
field makes the lifted difference principal. Its coefficient vector
therefore lies in the full geometric cuspidal modular-unit lattice
and is fixed by Frobenius $5$.

We compute this lattice from the integral period map for
$\Gamma_\Delta(72)$, where $\Delta=\{\pm1,\pm35\}$.
There are $64$ geometric cusps and $36$ Frobenius orbits.
Intersecting the complete principal-divisor lattice with the
Frobenius-fixed cusp-divisor lattice gives a rank-$35$ lattice
in $\mathbb Z^{64}$. Ordinary Euclidean norm in these coordinates
is exactly the weighted norm above.

Complete short-vector enumeration gives $24552$ pairs of nonzero
vectors of squared norm at most $104$. None has positive part of
degree $10$. An independent calculation reconstructs the full
period kernel and obtains the same enumeration using PARI.
This contradicts the vector $A_1-A_2$.
\end{proof}


\begin{proposition}\label{prop:decic236-excluded}
There is no effective divisor of degree $10$ on $X_1(2,36)$,
defined over $\mathbb Q$, whose support is entirely noncuspidal.
Moreover, $X_1(2,36)$ has rational gonality at least $11$.
Consequently neither $C_2\times C_{36}$ nor $C_{72}$ occurs as
a torsion subgroup over a number field of degree $10$.
\end{proposition}
\begin{proof}
Lemmas~\ref{lem:decic236-finite-sieve} and
\ref{lem:236-degree10-oldN8} exclude every such divisor whose
reduction has a noncuspidal component.

By \cite[Table~1, the entry for $X_1(2,36)$ at $5$]{NV},
\[
 \operatorname{gon}_{\mathbb F_5}(X_{\mathbb F_5})\geq10.
\]
Suppose that $X$ has a rational pencil of degree at most $10$.
The finite-field bound implies that its degree is $10$.
Take a saturated integral lattice in its two-dimensional section
space. Its reduction gives two independent sections of a
line bundle of degree $10$.
A common base divisor would give a morphism of smaller degree
over $\mathbb F_5$, contradicting the same bound.
Thus the reduced pencil is a degree-$10$ morphism.
By Lemma~\ref{lem:decic236-cuspidal-pencil}, it has a rational
fibre containing a noncuspidal point.
Choose a rational parameter reducing to that fibre and avoiding
the finitely many images of global cusps.
The resulting rational effective divisor has degree $10$,
entirely noncuspidal support, and a noncuspidal component in its
reduction. This is impossible by the preceding paragraph.
Hence the rational gonality is at least $11$.

Now let $D$ be an effective rational divisor of degree $10$
with entirely noncuspidal support and entirely cuspidal reduction.
Lift its reduction to a cuspidal divisor $C$ over a finite
unramified extension of $\mathbb Q_5$.
Choose a rational cusp $c$ of $X$.
The class $[D-10c]$ is torsion because $J_X(\mathbb Q)$ has
rank zero, and $[C-10c]$ is torsion by Manin--Drinfeld.
Thus $[D-C]$ is torsion and reduces to zero.
Torsion reduction is injective over an unramified extension of
$\mathbb Q_5$, so $D\sim C$.
These are distinct effective divisors. It follows that
$h^0(X,\mathcal O(D))\geq2$ after extending the base field,
and hence also over $\mathbb Q$.
This contradicts the rational gonality bound.

A point of $X$ over a degree-$10$ field gives, by its trace, an
effective rational divisor of degree $10$ with noncuspidal
support. This excludes $C_2\times C_{36}$.
Finally, suppose that $R\in E(K)$ has order $72$.
Let $\phi:E\to E'=E/\langle36R\rangle$.
The point $Q=\phi(R)$ has order $36$.
The nonzero point $P$ in the kernel of the dual isogeny is
$K$-rational and has order $2$.
It differs from $18Q$, since the dual isogeny sends $18Q$ to
$36R\ne0$.
Thus $E'(K)$ contains $C_2\times C_{36}$, a contradiction.
\end{proof}
